\documentclass[final]{beamer}

% ============================================================================
% 1. FORMATO GENERAL
% ============================================================================
% beamerposter permite conservar el requisito de la Escuela: A0 vertical.
% IEEEtran no se usa como clase porque esta diseñada para articulos, no para
% posters. Las convenciones IEEE se aplican a la estructura y las referencias.
\usepackage[orientation=portrait,size=a0,scale=1.0]{beamerposter}
\usepackage[T1]{fontenc}
\usepackage[utf8]{inputenc}
\usepackage[spanish,es-tabla]{babel}
\usepackage{newtxtext,newtxmath}
\usepackage{microtype}
\usepackage{ragged2e}
\usepackage{booktabs}
\usepackage{tabularx}
\usepackage{array}
\usepackage{tikz}
\usetikzlibrary{arrows.meta,positioning,calc}
\usepackage{url}
\usepackage{cite}
\usepackage{graphicx}

\setbeamertemplate{navigation symbols}{}
\setbeamertemplate{bibliography item}{\insertbiblabel}
\urlstyle{same}

% ============================================================================
% 2. PALETA Y TIPOGRAFIA
% ============================================================================
\definecolor{Navy}{HTML}{000000}
\definecolor{IEEEBlue}{HTML}{000000}
\definecolor{BrightBlue}{HTML}{000000}
\definecolor{Teal}{HTML}{000000}
\definecolor{PaleBlue}{HTML}{FFFFFF}
\definecolor{PaleTeal}{HTML}{FFFFFF}
\definecolor{PaleGray}{HTML}{F2F2F2}
\definecolor{Coral}{HTML}{000000}
\definecolor{BodyText}{HTML}{000000}
\definecolor{RuleGray}{HTML}{777777}
\definecolor{White}{HTML}{FFFFFF}

\setbeamercolor{normal text}{fg=BodyText,bg=White}
\setbeamercolor{headline}{fg=Navy,bg=White}
\setbeamercolor{block title}{fg=BodyText,bg=White}
\setbeamercolor{block body}{fg=BodyText,bg=White}
\setbeamercolor{alerted text}{fg=Coral}
\setbeamercolor{trace}{fg=Navy,bg=White}
\setbeamercolor{finding number}{fg=White,bg=Coral}
\setbeamercolor{conclusion}{fg=BodyText,bg=White}
\setbeamercolor{repository}{fg=BodyText,bg=White}
\setbeamercolor{note}{fg=BodyText,bg=White}
\setbeamercolor{structure}{fg=BodyText}
\AtBeginDocument{\hypersetup{allcolors=black}}

\setbeamerfont{normal text}{size=\fontsize{31}{35}\selectfont}
\setbeamerfont{block title}{series=\mdseries,shape=\scshape,size=\fontsize{35}{39}\selectfont}
\setbeamerfont{block body}{size=\fontsize{31}{35}\selectfont}
\setlength{\parindent}{1.15em}
\setlength{\parskip}{0.08em}

% Encabezados de seccion centrados y sin relleno, como en el articulo IEEE
% usado como referencia visual.
\setbeamertemplate{block begin}{
  \vskip0.55cm
  \begin{beamercolorbox}[
    wd=\linewidth,
    sep=0.08cm,
    center
  ]{block title}
    \usebeamerfont{block title}\MakeUppercase{\insertblocktitle}
  \end{beamercolorbox}
  \vspace{0.28cm}
  \begin{beamercolorbox}[
    wd=\linewidth,
    sep=0cm
  ]{block body}
    \usebeamerfont{block body}\justifying
}
\setbeamertemplate{block end}{
  \end{beamercolorbox}
  \vskip0.25cm
}

% ============================================================================
% 3. COMANDOS REUTILIZABLES
% ============================================================================
\newcommand{\metric}[3][BrightBlue]{%
  \begin{minipage}[t]{0.47\linewidth}
    {\color{#1}\bfseries\fontsize{48}{51}\selectfont #2\par}
    \vspace{0.15cm}
    {\bfseries\fontsize{24}{28}\selectfont #3\par}
  \end{minipage}%
}

\newcommand{\finding}[3]{%
  \noindent
  {\bfseries\itshape #1) #2.}\enspace
  {\fontsize{27}{31}\selectfont #3\par}
  \vspace{0.40cm}
}

\newcommand{\smalltext}{\fontsize{25}{29}\selectfont}
\newcommand{\bodytext}{\fontsize{31}{35}\selectfont}

% ============================================================================
% 4. DATOS DEL POSTER
% ============================================================================
\title{Diseño arquitectónico mínimo para sistemas RAG observables en español latinoamericano}
\author{Nicole Tatiana Parra Valverde}
\institute{Trabajo independiente}
\date{Segunda Escuela Sudamericana de NLP -- Sesión 2, orden 12}

\begin{document}
\begin{frame}[t]

% ============================================================================
% 5. ENCABEZADO
% ============================================================================
\begin{beamercolorbox}[wd=\textwidth,sep=0cm]{headline}
  \centering
  {\color{BodyText}\scshape\fontsize{20}{24}\selectfont
    Segunda Escuela Sudamericana de NLP -- Buenos Aires -- 2026\par}
  \vspace{0.8cm}
  {\color{BodyText}\mdseries\fontsize{74}{80}\selectfont
    Diseño arquitectónico mínimo para sistemas RAG observables en español latinoamericano\par}
  \vspace{0.9cm}
  {\color{BodyText}\fontsize{29}{34}\selectfont
    Nicole Tatiana Parra Valverde\par}
  \vspace{0.18cm}
  {\color{BodyText}\itshape\fontsize{23}{28}\selectfont
    Trabajo independiente\par}
  \vspace{0.18cm}
  {\color{BodyText}\fontsize{21}{25}\selectfont
    Caso de estudio: normativa estudiantil del Tecnológico de Costa Rica\par}
\end{beamercolorbox}

\vspace{0.55cm}
\noindent{\color{BodyText}\rule{\textwidth}{0.045cm}}
\vspace{0.35cm}

% ============================================================================
% 6. ARQUITECTURA CENTRAL
% ============================================================================
\begin{beamercolorbox}[wd=\textwidth,sep=0cm]{block body}
  \centering
  \resizebox{0.985\linewidth}{!}{%
  \begin{tikzpicture}[
    stage/.style={
      draw=BodyText,
      line width=1.1pt,
      fill=White,
      minimum width=9.0cm,
      minimum height=5.3cm,
      text width=7.9cm,
      align=left,
      inner sep=0.55cm
    },
    stagegreen/.style={
      stage
    },
    stagealert/.style={
      stage
    },
    flow/.style={
      -{Stealth[length=0.55cm,width=0.42cm]},
      draw=BodyText,
      line width=1.5pt
    }
  ]
    \node[stage] (source) {
      {\bfseries 01}\\[0.25cm]
      {\bfseries\fontsize{24}{27}\selectfont Fuentes}\\[0.20cm]
      {\fontsize{18}{22}\selectfont 25 documentos\\captura + hash}
    };
    \node[stage,right=1.0cm of source] (sections) {
      {\bfseries 02}\\[0.25cm]
      {\bfseries\fontsize{24}{27}\selectfont Secciones}\\[0.20cm]
      {\fontsize{18}{22}\selectfont 1\,184 registros\\artículos + transitorios}
    };
    \node[stagegreen,right=1.0cm of sections] (index) {
      {\bfseries 03}\\[0.25cm]
      {\bfseries\fontsize{24}{27}\selectfont Índice}\\[0.20cm]
      {\fontsize{18}{22}\selectfont 1\,259 fragmentos\\E5 · 384 dimensiones}
    };
    \node[stagegreen,right=1.0cm of index] (retrieval) {
      {\bfseries 04}\\[0.25cm]
      {\bfseries\fontsize{24}{27}\selectfont Recuperación}\\[0.20cm]
      {\fontsize{18}{22}\selectfont top-5\\similitud coseno}
    };
    \node[stage,right=1.0cm of retrieval] (generation) {
      {\bfseries 05}\\[0.25cm]
      {\bfseries\fontsize{24}{27}\selectfont Generación}\\[0.20cm]
      {\fontsize{18}{22}\selectfont Qwen3.5:9B\\Ollama local}
    };
    \node[stagealert,right=1.0cm of generation] (validation) {
      {\bfseries 06}\\[0.25cm]
      {\bfseries\fontsize{24}{27}\selectfont Validación}\\[0.20cm]
      {\fontsize{18}{22}\selectfont citas visibles\\+ estructuradas}
    };
    \node[stagegreen,right=1.0cm of validation] (output) {
      {\bfseries 07}\\[0.25cm]
      {\bfseries\fontsize{24}{27}\selectfont Salida}\\[0.20cm]
      {\fontsize{18}{22}\selectfont respuesta\\+ traza JSON}
    };

    \draw[flow] (source.east) -- (sections.west);
    \draw[flow] (sections.east) -- (index.west);
    \draw[flow] (index.east) -- (retrieval.west);
    \draw[flow] (retrieval.east) -- (generation.west);
    \draw[flow] (generation.east) -- (validation.west);
    \draw[flow] (validation.east) -- (output.west);
  \end{tikzpicture}%
  }

  \vspace{0.35cm}
  {\fontsize{19}{23}\selectfont
    \textbf{Fig. 1.} Arquitectura RAG mínima y observable. Cada componente
    produce un artefacto verificable; la traza conserva pregunta, fragmentos,
    contexto, respuesta, estado, citas y tiempos.\par}

  \vspace{0.2cm}
  {\fontsize{18}{22}\selectfont
    \textit{Configuración:} multilingual-e5-small; top-5; Qwen3.5:9B
    Q4\_K\_M; temperatura 0,1; ejecución local.\par}
\end{beamercolorbox}

\vspace{0.35cm}

% ============================================================================
% 7. CUERPO EN DOS COLUMNAS
% ============================================================================
\begin{columns}[T,totalwidth=\textwidth]

% --------------------------------------------------------------------------
% COLUMNA 1: RESUMEN, METODO Y RESULTADOS PRINCIPALES
% --------------------------------------------------------------------------
\begin{column}{0.485\textwidth}

\begin{block}{I. Resumen y objetivo}
\noindent\textbf{\textit{Resumen—}}Los prototipos RAG suelen mostrar una
respuesta final, pero no siempre
conservan la cadena de evidencia necesaria para explicar por qué fue generada.
Sin una traza es difícil distinguir fallos de recuperación, generación o
validación. El diseño parte del paradigma RAG descrito por Lewis et al.
\cite{lewis2020rag}.

\vspace{0.3cm}
\noindent\textbf{\textit{Objetivo—}}Diseñar e implementar una arquitectura
RAG mínima, local y reproducible para documentos en español latinoamericano,
con señales básicas de observabilidad por consulta.

\vspace{0.3cm}
{\smalltext\noindent\textbf{\textit{Palabras clave—}}RAG, PLN en español, recuperación de
información, observabilidad, modelos de lenguaje.}
\end{block}

\begin{block}{II. Caso de estudio}
\begin{columns}[T,totalwidth=\linewidth]
  \begin{column}{0.30\linewidth}
    {\color{Teal}\bfseries\fontsize{42}{45}\selectfont 25\par}
    {\bfseries\smalltext documentos oficiales}
  \end{column}
  \begin{column}{0.34\linewidth}
    {\color{Teal}\bfseries\fontsize{42}{45}\selectfont 1\,184\par}
    {\bfseries\smalltext secciones procesadas}
  \end{column}
  \begin{column}{0.32\linewidth}
    {\color{Teal}\bfseries\fontsize{42}{45}\selectfont 1\,259\par}
    {\bfseries\smalltext fragmentos indexados}
  \end{column}
\end{columns}

\vspace{0.55cm}
Normativa estudiantil del TEC, Costa Rica. Cada captura conserva URL, fecha y
suma de verificación. El corpus, sus metadatos y el procesamiento están
versionados en el repositorio \cite{tec2026reglamentos}.
\end{block}

\begin{block}{III. Metodología}
\begin{enumerate}
  \item Recolección de 25 documentos oficiales.
  \item Extracción de artículos y disposiciones transitorias.
  \item Fragmentación con máximo de 480 tokens y solapamiento de 50.
  \item Embeddings con \texttt{multilingual-e5-small}
    \cite{wang2024e5}, normalizados a 384 dimensiones.
  \item Recuperación top-5 por similitud coseno.
  \item Generación local con Qwen3.5:9B y temperatura 0,1.
  \item Validación de estado, citas y contrato JSON.
\end{enumerate}

\vspace{0.25cm}
\textbf{Evaluación:} 40 preguntas de referencia: 25 directas, 8 de varias
secciones y 7 no respondibles. Cada ejecución guardó fragmentos, respuesta,
citas y tiempos.
\end{block}

\begin{block}{IV. Revisión humana}
Las 40 respuestas se revisaron manualmente con seis criterios:
\begin{itemize}
  \item corrección factual;
  \item cobertura;
  \item fidelidad a fuentes;
  \item claridad;
  \item presencia de evidencia gold;
  \item conclusión manual.
\end{itemize}

\vspace{0.35cm}
\noindent{\bfseries\smalltext
Las métricas automáticas describen recuperación, citas y latencia; no se
presentan como exactitud factual.}
\end{block}

\begin{block}{V. Resultados}
La evaluación humana permite distinguir contenido correcto, cobertura de
evidencia y validez de la salida.

\vspace{0.55cm}
\metric{78,8\,\%}{corrección factual estricta}
\hfill
\metric[Teal]{93,5\,\%}{fidelidad total a fuentes}

\vspace{0.65cm}
{\color{RuleGray}\rule{\linewidth}{0.04cm}}
\vspace{0.55cm}

\metric{72,7\,\%}{cobertura completa}
\hfill
\metric[Teal]{7 / 7}{abstenciones correctas}
\end{block}

\begin{block}{VI. Aprobación por tipo de pregunta}
\centering
\begin{tikzpicture}[x=0.065\linewidth,y=1.72cm]
  \draw[draw=RuleGray,line width=1.2pt] (0,0) -- (10,0);
  \foreach \x/\label in {0/0,2/20,4/40,6/60,8/80,10/100} {
    \draw[draw=RuleGray] (\x,0.05) -- (\x,-0.08);
    \node[below,font=\smalltext,text=BodyText] at (\x,-0.08) {\label\%};
  }

  \node[left,align=right,text width=6.0cm,font=\smalltext,text=BodyText]
    at (0,1.2) {Directas · 19/25};
  \fill[BrightBlue] (0,0.82) rectangle (7.6,1.52);
  \node[right,font=\bfseries\smalltext,text=Navy] at (7.6,1.17) {76\%};

  \node[left,align=right,text width=6.0cm,font=\smalltext,text=BodyText]
    at (0,2.35) {Varias secciones · 2/8};
  \fill[BrightBlue] (0,1.97) rectangle (2.5,2.67);
  \node[right,font=\bfseries\smalltext,text=Navy] at (2.5,2.32) {25\%};

  \node[left,align=right,text width=6.0cm,font=\smalltext,text=BodyText]
    at (0,3.5) {No respondibles · 7/7};
  \fill[BrightBlue] (0,3.12) rectangle (10,3.82);
  \node[left,font=\bfseries\smalltext,text=White] at (9.85,3.47) {100\%};
\end{tikzpicture}

\vspace{0.5cm}
\noindent\textbf{\textit{Hallazgo central—}}El sistema se abstiene bien, pero
pierde cobertura cuando debe combinar evidencia distribuida.
\end{block}

\end{column}

% --------------------------------------------------------------------------
% COLUMNA 2: RENDIMIENTO, DISCUSION Y CONCLUSION
% --------------------------------------------------------------------------
\begin{column}{0.485\textwidth}

\begin{block}{VII. Rendimiento automático}
\begin{columns}[T,totalwidth=\linewidth]
  \begin{column}{0.48\linewidth}
    {\color{IEEEBlue}\bfseries\fontsize{29}{32}\selectfont 93,9\,\%\par}
    {\bfseries recuperación con evidencia gold en top-5}
  \end{column}
  \begin{column}{0.46\linewidth}
    {\color{IEEEBlue}\bfseries\fontsize{29}{32}\selectfont 27 ms\par}
    {\bfseries búsqueda semántica promedio}
  \end{column}
\end{columns}

\vspace{0.65cm}
\noindent\textbf{5,53 s por pregunta:} la generación consume 5,36 s en
promedio y domina la latencia del pipeline. La búsqueda representa menos del
1\,\% del tiempo total.
\end{block}

\begin{block}{VIII. Interpretación de las métricas}
\begin{itemize}
  \item \textbf{Corrección factual estricta:} proporción de respuestas
    respondibles calificadas como correctas, sin dar crédito a las parciales.
  \item \textbf{Fidelidad total:} grado en que las afirmaciones permanecen
    apoyadas por los fragmentos recuperados.
  \item \textbf{Cobertura completa:} presencia de toda la evidencia necesaria,
    no solo de una parte.
  \item \textbf{Recuperación top-5:} localización automática de al menos una
    evidencia gold entre los cinco primeros fragmentos.
\end{itemize}
\end{block}

\begin{block}{IX. Lo que reveló la observabilidad}
\finding{1}{Evidencia distribuida}{
  Solo 2 de 8 preguntas de varias secciones fueron aprobadas. Con frecuencia
  se recuperó una parte, pero no todos los artículos necesarios.}

\finding{2}{Documento correcto, fragmento incorrecto}{
  En q029 apareció el artículo relacionado, pero no el fragmento que contenía
  el dato. El análisis por identificador de registro habría ocultado el fallo.}

\finding{3}{Contenido correcto, salida inválida}{
  q025 y q032 se rechazaron porque las citas visibles y estructuradas no
  coincidían. La validez del contrato también es calidad.}

\finding{4}{Abstención confiable}{
  Las 7 preguntas sobre precios, horarios, contraseñas o cupos actuales
  recibieron \texttt{not\_found}, sin inventar información ausente.}
\end{block}

\begin{block}{X. Conclusión}
\noindent{\bfseries
Una arquitectura pequeña puede ofrecer trazabilidad útil sin una plataforma
compleja. Registrar evidencia, citas, estados y tiempos permitió localizar
fallos en recuperación, generación y validación.}

\vspace{0.35cm}
El prototipo funcionó mejor con preguntas directas y abstenciones. El reto
principal fue combinar evidencia distribuida. Por ello, una arquitectura RAG
observable debe distinguir evidencia por documento y por fragmento, y validar
la consistencia de las citas antes de entregar la respuesta.
\end{block}

\begin{block}{XI. Limitaciones y trabajo futuro}
\textbf{Limitaciones}
\begin{itemize}
  \item un dominio, un embedding y un modelo generativo;
  \item 40 preguntas y una persona revisora;
  \item sin comparación léxica, híbrida ni reranking;
  \item el caso no representa toda la diversidad del español latinoamericano.
\end{itemize}

\vspace{0.35cm}
\textbf{Trabajo futuro}
\begin{itemize}
  \item reranking para evidencia multisección;
  \item comparar recuperación densa, léxica e híbrida;
  \item ampliar países, dominios y revisión humana;
  \item estudiar umbrales de abstención y calibración.
\end{itemize}
\end{block}

\begin{block}{Código, datos y documentación}
\centering
{\smalltext Repositorio reproducible:\par}
\vspace{0.15cm}
{\bfseries\fontsize{25}{29}\selectfont
\url{https://github.com/niparra21/nlpArgentina}}
\end{block}

\end{column}
\end{columns}

\vspace{0.45cm}

% ============================================================================
% 8. REFERENCIAS IEEE Y PIE
% ============================================================================
\begin{beamercolorbox}[wd=\textwidth,sep=0cm]{block body}
  {\color{Navy}\bfseries\fontsize{26}{30}\selectfont XII. Referencias\par}
  \vspace{0.25cm}
  {\fontsize{16}{20}\selectfont
    \bibliographystyle{IEEEtran}
    \bibliography{references}
  }
\end{beamercolorbox}

\vspace{0.25cm}
\begin{columns}[T,totalwidth=\textwidth]
  \begin{column}{0.76\textwidth}
    {\color{BodyText}\fontsize{15}{18}\selectfont
      Los resultados corresponden a una corrida reproducible del prototipo y
      no sustituyen la consulta de la normativa oficial.}
  \end{column}
  \begin{column}{0.22\textwidth}
    \raggedleft
    {\color{Navy}\bfseries\fontsize{18}{21}\selectfont Sesión 2 · Orden 12}
  \end{column}
\end{columns}

\end{frame}
\end{document}
